\begin{pastexam}[label={px:df3-107-ds-03}]{地方三等資訊處理 107 年 第三題}
（申論；20 分）

假設陣列 $A[1..n]$ 儲存 $n$ 個正整數 $x1, x2,\ldots, x_n$。（每小題 10 分，共 20 分）
\begin{itemize}[leftmargin=2.2em]
\item[㈠] 已知所有的正整數 $x_i \le M$。請設計一個 $O(n+M)$ 時間的演算法將這些整數由小到大排列。
\item[㈡] 已知所有的正整數 $x_i \le n^2$。請設計一個 $O(n)$ 時間的演算法將這些整數由小到大排列，或證明這是不可行的。
\end{itemize}
\end{pastexam}

\begin{pastsoln}
原卷的「$x1, x2$」是 $x_1, x_2$ 漏了下標。 這兩小題都是不靠比較的線性時間排序（資結書第~15.2~節）。 比較排序有
$\Omega(n \log n)$ 的下界（資結書定理~15.1），但這裡知道鍵值的範圍，可以繞過下界。

\textbf{㈠ 計數排序}（資結書虛擬碼~15.1）。 開一個計數陣列 $C[1..M]$，先數每個值出現幾次，再算前綴和得到每個值在輸出中的
最後位置，最後由後往前把元素放到輸出陣列 $B$：
\begin{pseudo}[unbreakable]{計數排序，$1 \le A[i] \le M$}
\Procedure{Counting-Sort}{$A, B, n, M$}
    \For{$k \gets 1$ \To $M$}
        \State $C[k] \gets 0$
    \EndFor
    \For{$i \gets 1$ \To $n$}
        \State $C[A[i]] \gets C[A[i]] + 1$ \Comment{$C[k]$ = 值為 $k$ 的個數}
    \EndFor
    \For{$k \gets 2$ \To $M$}
        \State $C[k] \gets C[k] + C[k - 1]$ \Comment{$C[k]$ = 值 $\le k$ 的個數}
    \EndFor
    \For{$i \gets n$ \DownTo $1$}
        \State $B[C[A[i]]] \gets A[i]$；$C[A[i]] \gets C[A[i]] - 1$
    \EndFor
\EndProcedure
\end{pseudo}
\textbf{正確性}：第 6、7 行做完後，$C[k]$ 是 $\le k$ 的元素個數，所以值為 $k$ 的元素在輸出中應該佔位置 $C[k-1]+1..C[k]$；第 8、9 行由後往前放，
每放一個就把 $C[A[i]]$ 減一，值為 $k$ 的元素依序填進這一段，而且相同值保持原本的相對順序（穩定）。 \textbf{時間}：四個迴圈分別是
$O(M)$、$O(n)$、$O(M)$、$O(n)$，合計 $O(n + M)$；額外空間 $O(n + M)$。 若只要排序整數本身，也可以數完以後依 $k = 1..M$ 把 $k$ 寫出
$C[k]$ 次，同樣 $O(n + M)$。

\textbf{㈡ 可行：以 $n$ 為基數的基數排序}。 直接用 ㈠ 取 $M = n^2$ 是 $O(n^2)$，不夠快。 改成把每個數看成\textbf{以 $n$ 為底的兩位數}：
令 $y_i = x_i - 1$，則 $0 \le y_i \le n^2 - 1$，可唯一寫成
\[
y_i = q_i \cdot n + r_i,\qquad q_i = \lfloor y_i / n \rfloor,\quad r_i = y_i \bmod n,\qquad 0 \le q_i, r_i \le n - 1.
\]
先用計數排序依\textbf{低位} $r_i$ 排（穩定），再用計數排序依\textbf{高位} $q_i$ 排（穩定），即 LSD 基數排序（資結書第~15.2.2~節，虛擬碼~15.2）。
每一位的範圍是 $0..n-1$，一次計數排序 $O(n + n) = O(n)$；兩位共 $O(n)$。 （計算 $q_i$、$r_i$ 每個 $O(1)$，共 $O(n)$。）

\textbf{正確性}：兩個數 $y_i < y_j$ 時，若 $q_i < q_j$，第二趟依高位排就把 $y_i$ 放前面；若 $q_i = q_j$，則 $r_i < r_j$，第一趟後 $y_i$ 在
$y_j$ 前面，第二趟高位相同，穩定排序保持這個順序。 所以兩趟後依 $y$（也就是 $x$）由小到大。 $x_i \le n^2$ 的條件保證只要兩位；
一般來說 $x_i \le n^k$ 用 $k$ 趟，$O(kn)$，$k$ 是常數時仍是線性。

例：$n = 4$、$x = 13, 2, 16, 7$（都 $\le 16$），則 $y = 12, 1, 15, 6$，寫成 $(q, r)$ 是 $(3,0)$、$(0,1)$、$(3,3)$、$(1,2)$。
依 $r$ 排：$(3,0)$、$(0,1)$、$(1,2)$、$(3,3)$；再依 $q$ 穩定排：$(0,1)$、$(1,2)$、$(3,0)$、$(3,3)$，即 $y = 1, 6, 12, 15$，
排好的 $x$ 是 $2, 7, 13, 16$。
\end{pastsoln}
